\documentclass[11pt,preprint]{elsarticle}

\usepackage{lmodern}
%%%% My spacing
\usepackage{setspace}
\setstretch{1.2}
\DeclareMathSizes{12}{14}{10}{10}

% Wrap around which gives all figures included the [H] command, or places it "here". This can be tedious to code in Rmarkdown.
\usepackage{float}
\let\origfigure\figure
\let\endorigfigure\endfigure
\renewenvironment{figure}[1][2] {
    \expandafter\origfigure\expandafter[H]
} {
    \endorigfigure
}

\let\origtable\table
\let\endorigtable\endtable
\renewenvironment{table}[1][2] {
    \expandafter\origtable\expandafter[H]
} {
    \endorigtable
}


\usepackage{ifxetex,ifluatex}
\usepackage{fixltx2e} % provides \textsubscript
\ifnum 0\ifxetex 1\fi\ifluatex 1\fi=0 % if pdftex
  \usepackage[T1]{fontenc}
  \usepackage[utf8]{inputenc}
\else % if luatex or xelatex
  \ifxetex
    \usepackage{mathspec}
    \usepackage{xltxtra,xunicode}
  \else
    \usepackage{fontspec}
  \fi
  \defaultfontfeatures{Mapping=tex-text,Scale=MatchLowercase}
  \newcommand{\euro}{€}
\fi

\usepackage{amssymb, amsmath, amsthm, amsfonts}

\def\bibsection{\section*{References}} %%% Make "References" appear before bibliography


\usepackage[numbers]{natbib}

\usepackage{longtable}
\usepackage[margin=2.3cm,bottom=2cm,top=2.5cm, includefoot]{geometry}
\usepackage{fancyhdr}
\usepackage[bottom, hang, flushmargin]{footmisc}
\usepackage{graphicx}
\numberwithin{equation}{section}
\numberwithin{figure}{section}
\numberwithin{table}{section}
\setlength{\parindent}{0cm}
\setlength{\parskip}{1.3ex plus 0.5ex minus 0.3ex}
\usepackage{textcomp}
\renewcommand{\headrulewidth}{0.2pt}
\renewcommand{\footrulewidth}{0.3pt}

\usepackage{array}
\newcolumntype{x}[1]{>{\centering\arraybackslash\hspace{0pt}}p{#1}}

%%%%  Remove the "preprint submitted to" part. Don't worry about this either, it just looks better without it:
\makeatletter
\def\ps@pprintTitle{%
  \let\@oddhead\@empty
  \let\@evenhead\@empty
  \let\@oddfoot\@empty
  \let\@evenfoot\@oddfoot
}
\makeatother

 \def\tightlist{} % This allows for subbullets!

\usepackage{hyperref}
\hypersetup{breaklinks=true,
            bookmarks=true,
            colorlinks=true,
            citecolor=blue,
            urlcolor=blue,
            linkcolor=blue,
            pdfborder={0 0 0}}


% The following packages allow huxtable to work:
\usepackage{siunitx}
\usepackage{multirow}
\usepackage{hhline}
\usepackage{calc}
\usepackage{tabularx}
\usepackage{booktabs}
\usepackage{caption}


\newenvironment{columns}[1][]{}{}

\newenvironment{column}[1]{\begin{minipage}{#1}\ignorespaces}{%
\end{minipage}
\ifhmode\unskip\fi
\aftergroup\useignorespacesandallpars}

\def\useignorespacesandallpars#1\ignorespaces\fi{%
#1\fi\ignorespacesandallpars}

\makeatletter
\def\ignorespacesandallpars{%
  \@ifnextchar\par
    {\expandafter\ignorespacesandallpars\@gobble}%
    {}%
}
\makeatother


% definitions for citeproc citations
\NewDocumentCommand\citeproctext{}{}
\NewDocumentCommand\citeproc{mm}{%
\href{\#cite.\detokenize{#1}}{#2}\nocite{#1}}

\makeatletter
% allow citations to break across lines
\let\@cite@ofmt\@firstofone
% avoid brackets around text for \cite:
\def\@biblabel#1{}
\def\@cite#1#2{{#1\if@tempswa , #2\fi}}
\makeatother
\newlength{\cslhangindent}
\setlength{\cslhangindent}{1.5em}
\newlength{\csllabelwidth}
\setlength{\csllabelwidth}{3em}
\newenvironment{CSLReferences}[2] % #1 hanging-indent, #2 entry-spacing
{\begin{list}{}{%
	\setlength{\itemindent}{0pt}
	\setlength{\leftmargin}{0pt}
	\setlength{\parsep}{0pt}
	% turn on hanging indent if param 1 is 1
	\ifodd #1
	\setlength{\leftmargin}{\cslhangindent}
	\setlength{\itemindent}{-1\cslhangindent}
	\fi
	% set entry spacing
	\setlength{\itemsep}{#2\baselineskip}}}
{\end{list}}

\usepackage{calc}
\newcommand{\CSLBlock}[1]{\hfill\break\parbox[t]{\linewidth}{\strut\ignorespaces#1\strut}}
\newcommand{\CSLLeftMargin}[1]{\parbox[t]{\csllabelwidth}{\strut#1\strut}}
\newcommand{\CSLRightInline}[1]{\parbox[t]{\linewidth - \csllabelwidth}{\strut#1\strut}}
\newcommand{\CSLIndent}[1]{\hspace{\cslhangindent}#1}


\urlstyle{same}  % don't use monospace font for urls
\setlength{\parindent}{0pt}
\setlength{\parskip}{6pt plus 2pt minus 1pt}
\setlength{\emergencystretch}{3em}  % prevent overfull lines
\setcounter{secnumdepth}{5}

%%% Use protect on footnotes to avoid problems with footnotes in titles
\let\rmarkdownfootnote\footnote%
\def\footnote{\protect\rmarkdownfootnote}
\IfFileExists{upquote.sty}{\usepackage{upquote}}{}

%%% Include extra packages specified by user

%%% Hard setting column skips for reports - this ensures greater consistency and control over the length settings in the document.
%% page layout
%% paragraphs
\setlength{\baselineskip}{12pt plus 0pt minus 0pt}
\setlength{\parskip}{12pt plus 0pt minus 0pt}
\setlength{\parindent}{0pt plus 0pt minus 0pt}
%% floats
\setlength{\floatsep}{12pt plus 0 pt minus 0pt}
\setlength{\textfloatsep}{20pt plus 0pt minus 0pt}
\setlength{\intextsep}{14pt plus 0pt minus 0pt}
\setlength{\dbltextfloatsep}{20pt plus 0pt minus 0pt}
\setlength{\dblfloatsep}{14pt plus 0pt minus 0pt}
%% maths
\setlength{\abovedisplayskip}{12pt plus 0pt minus 0pt}
\setlength{\belowdisplayskip}{12pt plus 0pt minus 0pt}
%% lists
\setlength{\topsep}{10pt plus 0pt minus 0pt}
\setlength{\partopsep}{3pt plus 0pt minus 0pt}
\setlength{\itemsep}{5pt plus 0pt minus 0pt}
\setlength{\labelsep}{8mm plus 0mm minus 0mm}
\setlength{\parsep}{\the\parskip}
\setlength{\listparindent}{\the\parindent}
%% verbatim
\setlength{\fboxsep}{5pt plus 0pt minus 0pt}



\begin{document}



\begin{frontmatter}  %

\title{Question 4: Global Balanced Index Fund Construction}

% Set to FALSE if wanting to remove title (for submission)




\author[Add1]{24693197}
\ead{24693197@sun.ac.za}

\author[Add1,Add2]{John Smith}
\ead{John@gmail.com}

\author[Add1,Add2]{John Doe}
\ead{Joe@gmail.com}



\address[Add1]{Stellenbosch University, Stellenbosch, South Africa}
\address[Add2]{Some other Institution, Cape Town, South Africa}

\cortext[cor]{Corresponding author: 24693197}

\begin{abstract}
\small{
Abstract to be written here. The abstract should not be too long and
should provide the reader with a good understanding what you are writing
about. Academic papers are not like novels where you keep the reader in
suspense. To be effective in getting others to read your paper, be as
open and concise about your findings here as possible. Ideally, upon
reading your abstract, the reader should feel he / she must read your
paper in entirety.
}
\end{abstract}

\vspace{1cm}


\begin{keyword}
\footnotesize{
Multivariate GARCH \sep Kalman Filter \sep Copula \\
\vspace{0.3cm}
}
\footnotesize{
\textit{JEL classification} L250 \sep L100
}
\end{keyword}



\vspace{0.5cm}

\end{frontmatter}

\setcounter{footnote}{0}



%________________________
% Header and Footers
%%%%%%%%%%%%%%%%%%%%%%%%%%%%%%%%%
\pagestyle{fancy}
\chead{}
\rhead{}
\lfoot{}
\rfoot{\footnotesize Page \thepage}
\lhead{}
%\rfoot{\footnotesize Page \thepage } % "e.g. Page 2"
\cfoot{}

%\setlength\headheight{30pt}
%%%%%%%%%%%%%%%%%%%%%%%%%%%%%%%%%
%________________________

\headsep 35pt % So that header does not go over title




\section{\texorpdfstring{Introduction
\label{Introduction}}{Introduction }}\label{introduction}

Note: I was not sure what the desired output for this was so I just went
with a Texevier pdf :)

This report will construct a Global Balanced Index Fund using a mix of
traded global and local indexes. A balanced fund holds a mix of asset
classes, such as equities, bonds and commodities, and typically follows
a strategic asset allocation that is re-evaluated over time. As the fund
is built from indexes, the main decision is how much to allocate to each
index. The fund is rebalanced every quarter and must meet a set of
limits on its exposures, which are set out in Section
\ref{Construction}.

\section{\texorpdfstring{Data and approach
\label{Data}}{Data and approach }}\label{data-and-approach}

This section will look at the assets that the fund can hold:

\begin{itemize}
\tightlist
\item
  Global equity: MSCI USA, Europe, Japan and Emerging Markets indexes.
\item
  Local equity: the J200, the Mid Cap index (J201), the value and
  momentum style indexes and listed property (J253).
\item
  Bonds: the ALBI and inflation-linked bonds (CILI) locally, and US,
  German and UK 10 year government bonds globally.
\item
  Commodities: gold and the Bloomberg Commodity Index.
\end{itemize}

The bond data is given as yields, so these were first converted to
monthly returns (as per the instructions in the exam paper). The return
is calculated as last month's yield divided by twelve, minus the
modified duration multiplied by the change in yield. Modified duration
refers to how sensitive a bond's price is to a change in its yield and 9
was used based on the average yield since 2010 (see root readme).
Therefore, the bond returns are an approximation.

\begin{table}[H]
\centering
\begin{tabular}{llllr}
  \hline
Name & Class & Region & Start & Months \\ 
  \hline
Germany\_10Yr & Bonds & Global & Feb 2005 & 260 \\ 
  UK\_10Yr & Bonds & Global & Feb 2005 & 260 \\ 
  US\_10Yr & Bonds & Global & Feb 2005 & 260 \\ 
  ALBITR & Bonds & Local & Feb 2001 & 308 \\ 
  CILITR & Bonds & Local & Mar 2007 & 235 \\ 
  Bcom\_Index & Commodities & Global & Feb 2005 & 260 \\ 
  Gold & Commodities & Global & Feb 2005 & 260 \\ 
  MSCI\_EM & Equity & Global & Feb 2005 & 260 \\ 
  MSCI\_EU & Equity & Global & Feb 2005 & 260 \\ 
  MSCI\_Jap & Equity & Global & Feb 2005 & 260 \\ 
  MSCI\_USA & Equity & Global & Feb 2005 & 260 \\ 
  J200 & Equity & Local & Feb 2015 & 140 \\ 
  J201 & Equity & Local & Feb 2015 & 140 \\ 
  J253 & Equity & Local & Jan 2002 & 297 \\ 
  STXMOMJ433 & Equity & Local & Apr 2017 & 114 \\ 
  STXVALJ433 & Equity & Local & Apr 2018 & 102 \\ 
   \hline
\end{tabular}
\caption{Assets available to the fund \label{Table_Assets}} 
\end{table}

\section{\texorpdfstring{Constructing the fund
\label{Construction}}{Constructing the fund }}\label{constructing-the-fund}

This section will look at how the weights of the fund are chosen. The
fund must meet the following requirements:

\begin{itemize}
\tightlist
\item
  Long-only strategy.
\item
  When using covariance and mean forecasts, use a look-back of at least
  10 years.
\item
  Do not hold any assets with less than 5 years' returns data.
\item
  Apply quarterly rebalancing.
\item
  Limit exposure to bonds and credit instruments at 25\%.
\item
  Limit exposure to equities at 75\%.
\item
  Limit exposure to commodities to 15\%.
\item
  Limit single asset exposure to 35\%.
\item
  Global exposure limit 60\%.
\end{itemize}

The fund is also fully invested, tbus the weights sum to one.

With these limits kept in mind, the weights are chosen so that the fund
has the lowest possible variance. Minimum variance is an objective
optimisation technique, as it does not depend on a risk or return
preference. The accuracy of an optimiser out of sample depends on the
quality of the return forecasts and the stability of the covariance
matrix. Therefore, minimum variance was chosen as it does not need
return forecasts. It only uses the covariance matrix, and this shows how
the assets move together. Outliers have a large effect on the covariance
matrix. For this reason, it is estimated using Ledoit-Wolf shrinkage.
Shrinkage pulls the most extreme coefficients toward more central values
and reduces estimation error. The covariance matrix also needs a full
square matrix of returns with no missing values. Some assets only have
between five and ten years of data and therefore have gaps in the ten
year window. These gaps are filled with random draws from each asset's
own return distribution. This preserves the distribution of each asset.

\section{\texorpdfstring{Results
\label{Results}}{Results }}\label{results}

This section will look at how the fund performed compared to two other
portfolios. The first is a strategic balanced allocation, which holds
fixed weights based on the published benchmark weights of a South
African balanced index fund (see Satrix Balanced Index fund pdf in the
bin folder). These are 40\% in the Capped SWIX, 30\% in MSCI World, 8\%
in MSCI Emerging Markets, 17\% in the ALBI, 3\% in listed property and
2\% in cash (STeFI). Its weights are reset on the same quarterly dates
as the fund. The second is the Capped SWIX on its own, which shows how
the fund compares to holding only local equities.

\begin{figure}[H]

{\centering \includegraphics{Question4_files/figure-latex/Figure_Cum-1} 

}

\caption{Cumulative returns \label{Figure_Cum}}\label{fig:Figure_Cum}
\end{figure}

Figure \ref{Figure_Cum} shows the growth of R1 invested in the fund and
the two comparison portfolios since January 2020. As all three
portfolios start on the same date, the comparison is not distorted by
different start dates. The fund fell by less than both comparison
portfolios during the 2020 market decline (Covid period). However, it
grew more slowly from 2022 onwards and ends up below both the strategic
balanced allocation and the Capped SWIX at the end of the period. This
suggests that the fund offers more protection when markets fall, but
that clearly comes with the cost of gains when the market experiences an
upswing.

\begin{figure}[H]

{\centering \includegraphics{Question4_files/figure-latex/Figure_DD-1} 

}

\caption{Drawdowns \label{Figure_DD}}\label{fig:Figure_DD}
\end{figure}

Figure \ref{Figure_DD} shows the drawdowns of each portfolio, which
measure how far a portfolio falls from its previous peak. The fund had
the smallest drawdown in 2020 and near the end of the sample. In 2022,
however, its drawdown was similar to that of the strategic balanced
allocation. Therefore, the lower volatility of the fund did not always
lead to smaller losses.

\begin{figure}[H]

{\centering \includegraphics{Question4_files/figure-latex/Figure_Weights-1} 

}

\caption{Weights plot \label{Figure_Weights}}\label{fig:Figure_Weights}
\end{figure}

Figure \ref{Figure_Weights} shows the weights of the fund at each
rebalance. The limits on bonds, commodities and global assets are
binding at every rebalance, so the split between asset classes stays the
same throughout. The optimiser therefore mainly decides which assets to
hold within each class. Gold makes up most of the commodity weight for
most of the period. The local equity indices only enter the fund once
they have five years of data and MSCI Japan has the largest global
equity weight (over the period). As the global returns are in USD and
the local returns in ZAR, part of the preference for global assets may
be due to currency movements. Therefore, the weights should be
interpreted with caution.

Table \ref{Table_Stats} compares the three portfolios over the last
three and five years, with the strategic balanced allocation as the
benchmark. The fund has the lowest volatility over both periods, as well
as the smallest average drawdown and the smallest modified CVaR. Its
beta is also below one, which means that it moves less than the
benchmark. However, this lower risk came with lower returns (as we saw
in the cummulative returns figure). Over five years, the fund's
annualised return is below that of the benchmark and its information
ratio is negative over both periods. Over the last three years, the fund
earned a similar return to the Capped SWIX with much lower volatility.
This suggests that the fund achieves what a minimum variance approach
aims for (lower risk), but it did so at the cost of some of the returns.

\section{\texorpdfstring{Conclusion
\label{Conclusion}}{Conclusion }}\label{conclusion}

This report constructed a Global Balanced Index Fund that minimises
variance within the given limits and is rebalanced every quarter. The
limits on bonds, commodities and global assets were binding throughout.
Thus, the split between asset classes was largely set by the limits and
the optimiser mainly chose the assets within each class. The fund had
lower volatility than both the strategic balanced allocation and the
Capped SWIX and it declined by less during sharp equity declines.
However, it also earned lower returns and its losses were not always
smaller than those of the benchmark. These results should be interpreted
with caution, due to the short sample and the global returns are in USD
and the bond returns are approximations. As the output of an optimiser
is a guide rather than the truth, the weights are best used as a
starting point for further discussion.

\newpage

\section*{References}\label{references}
\addcontentsline{toc}{section}{References}

\protect\phantomsection\label{refs}

\section*{Appendix}\label{appendix}
\addcontentsline{toc}{section}{Appendix}

\subsection*{Appendix A}\label{appendix-a}
\addcontentsline{toc}{subsection}{Appendix A}

Some appendix information here

\subsection*{Appendix B}\label{appendix-b}
\addcontentsline{toc}{subsection}{Appendix B}

\bibliography{Tex/ref}





\end{document}
